\documentclass[10pt,a4paper]{amsart}
\usepackage{amsmath,amssymb,amsthm,graphicx}
% Retained preamble for the Stacks Project chapter snapshot in this directory.
% The chapter's own preamble.tex is a copy of the shared Stacks Project preamble; the block
% below is sliced MECHANICALLY, line for line, from it. Nothing was added, removed or
% redefined; the wrapper class, the xr-hyper/hyperref wiring and the externaldocument list
% are omitted because the assembler supplies the document class.
%
% lines taken from preamble.tex: 3, 5, 9, 11, 12, 156, 157, 158, ... (51 lines total)
%
% package coverage restored from the chapter's own preamble (2026-10-13): xy, xy-2cell

\usepackage[all]{xy}
\xyoption{2cell}
\UseAllTwocells

\newenvironment{reference}{\comment}{\endcomment}
\newenvironment{slogan}{\comment}{\endcomment}
\newenvironment{history}{\comment}{\endcomment}
\theoremstyle{plain}
\newtheorem{theorem}[subsection]{Theorem}
\newtheorem{proposition}[subsection]{Proposition}
\newtheorem{lemma}[subsection]{Lemma}
\theoremstyle{definition}
\newtheorem{definition}[subsection]{Definition}
\newtheorem{example}[subsection]{Example}
\newtheorem{exercise}[subsection]{Exercise}
\newtheorem{situation}[subsection]{Situation}
\theoremstyle{remark}
\newtheorem{remark}[subsection]{Remark}
\newtheorem{remarks}[subsection]{Remarks}
\numberwithin{equation}{subsection}
\def\lim{\mathop{\mathrm{lim}}\nolimits}
\def\colim{\mathop{\mathrm{colim}}\nolimits}
\def\Spec{\mathop{\mathrm{Spec}}}
\def\Hom{\mathop{\mathrm{Hom}}\nolimits}
\def\Ext{\mathop{\mathrm{Ext}}\nolimits}
\def\SheafHom{\mathop{\mathcal{H}\!\mathit{om}}\nolimits}
\def\SheafExt{\mathop{\mathcal{E}\!\mathit{xt}}\nolimits}
\def\Sch{\mathit{Sch}}
\def\Mor{\mathop{\mathrm{Mor}}\nolimits}
\def\Ob{\mathop{\mathrm{Ob}}\nolimits}
\def\Sh{\mathop{\mathit{Sh}}\nolimits}
\def\NL{\mathop{N\!L}\nolimits}
\def\CH{\mathop{\mathrm{CH}}\nolimits}
\def\proetale{{pro\text{-}\acute{e}tale}}
\def\etale{{\acute{e}tale}}
\def\QCoh{\mathit{QCoh}}
\def\Ker{\mathop{\mathrm{Ker}}}
\def\Im{\mathop{\mathrm{Im}}}
\def\Coker{\mathop{\mathrm{Coker}}}
\def\Coim{\mathop{\mathrm{Coim}}}
\DeclareMathSymbol{\boxtimes}{\mathbin}{AMSa}{"02}
\def\QCohstack{\mathcal{QC}\!\mathit{oh}}
\def\Cohstack{\mathcal{C}\!\mathit{oh}}
\def\Spacesstack{\mathcal{S}\!\mathit{paces}}
\def\Quotfunctor{\mathrm{Quot}}
\def\Hilbfunctor{\mathrm{Hilb}}
\def\Curvesstack{\mathcal{C}\!\mathit{urves}}
\def\Polarizedstack{\mathcal{P}\!\mathit{olarized}}
\def\Complexesstack{\mathcal{C}\!\mathit{omplexes}}
\def\Pic{\mathop{\mathrm{Pic}}\nolimits}
\def\Picardstack{\mathcal{P}\!\mathit{ic}}
\def\Picardfunctor{\mathrm{Pic}}
\def\Deformationcategory{\mathcal{D}\!\mathit{ef}}
\usepackage{cleveref}
\crefname{originallemma}{original lemma}{original lemmas}
\Crefname{originallemma}{Original lemma}{Original lemmas}
\makeatletter
\newlabel{definition-trace-element}{{0BT6}{source}{}{sourcepointer-1}{}}
\newlabel{definition-trace-element@cref}{{[originallemma][0][]0BT6}{[1][0][]source}{}{}{}}
\makeatother
\crefname{originallemma}{original lemma}{original lemmas}
\Crefname{originallemma}{Original lemma}{Original lemmas}
\makeatletter
\newlabel{equation-bc-dualizing}{{0BVB}{source}{}{sourcepointer-2}{}}
\newlabel{equation-bc-dualizing@cref}{{[originallemma][0][]0BVB}{[1][0][]source}{}{}{}}
\makeatother
\crefname{originallemma}{original lemma}{original lemmas}
\Crefname{originallemma}{Original lemma}{Original lemmas}
\makeatletter
\newlabel{lemma-dedekind-complementary-module}{{0BW2}{source}{}{sourcepointer-3}{}}
\newlabel{lemma-dedekind-complementary-module@cref}{{[originallemma][0][]0BW2}{[1][0][]source}{}{}{}}
\makeatother
\crefname{originallemma}{original lemma}{original lemmas}
\Crefname{originallemma}{Original lemma}{Original lemmas}
\makeatletter
\newlabel{lemma-dedekind-different-ideal}{{0BW1}{source}{}{sourcepointer-4}{}}
\newlabel{lemma-dedekind-different-ideal@cref}{{[originallemma][0][]0BW1}{[1][0][]source}{}{}{}}
\makeatother
\crefname{originallemma}{original lemma}{original lemmas}
\Crefname{originallemma}{Original lemma}{Original lemmas}
\makeatletter
\newlabel{lemma-dualizing-associated-primes}{{0BVD}{source}{}{sourcepointer-5}{}}
\newlabel{lemma-dualizing-associated-primes@cref}{{[originallemma][0][]0BVD}{[1][0][]source}{}{}{}}
\makeatother
\crefname{originallemma}{original lemma}{original lemmas}
\Crefname{originallemma}{Original lemma}{Original lemmas}
\makeatletter
\newlabel{lemma-explain-condition}{{0BWL}{source}{}{sourcepointer-6}{}}
\newlabel{lemma-explain-condition@cref}{{[originallemma][0][]0BWL}{[1][0][]source}{}{}{}}
\makeatother
\crefname{originallemma}{original lemma}{original lemmas}
\Crefname{originallemma}{Original lemma}{Original lemmas}
\makeatletter
\newlabel{remark-different-generalization}{{0BWM}{source}{}{sourcepointer-7}{}}
\newlabel{remark-different-generalization@cref}{{[originallemma][0][]0BWM}{[1][0][]source}{}{}{}}
\makeatother
\title{Stacks, Discriminants: lemma-agree-dedekind}
\author{The Stacks Project Authors}
\date{stacks-project:discriminant.tex:local-snapshot-20261012; GNU FDL 1.2 or later}
\begin{document}
\maketitle
\noindent\textit{Editorial note.} Statement span is lines 3268-3281 of sources/stacks-discriminant/pinned-original-source.tex; proof span is lines 3283-3321, the complete author proof (37 body lines).\par\medskip
\noindent\textit{Editorial note.} Navigation stubs: 7 cited-but-not-included labels are supplied as non-load-bearing pointers with their pinned tags where the snapshot tag index has them.\par\medskip
\noindent\textit{Editorial note.} extra\_packages is empty on purpose: the retained preamble already loads xy with options and a second declaration clashes in the assembler wrapper.\par\medskip
\noindent\textit{Editorial note.} Text is retained verbatim, including the author's own typos and placeholders.\par\medskip
\noindent\textit{Editorial note.} Queue index 13 of queue\_discriminant.json (pb=1585, stmt 3268-3281, proof 3283-3321).\par\medskip
\noindent\textit{Allowed standard prerequisites (editorial).}\begin{itemize}
\item Explicit prerequisite boundary. The cited inputs are prerequisites; the argument in the retained span is not replaced by any cited result.
\end{itemize}\medskip
\ifdefined\hypertarget\hypertarget{sourcepointer-1}{}\fi
\noindent\textit{Source reference.} Original lemma 0BT6: sources/stacks-discriminant/pinned-original-source.tex, source lines 821--836. Cited by the retained statement or proof but not part of the collected spans; supplied only as a navigation pointer. The argument does not depend on its content.\par\medskip
\ifdefined\hypertarget\hypertarget{sourcepointer-2}{}\fi
\noindent\textit{Source reference.} Original lemma 0BVB: sources/stacks-discriminant/pinned-original-source.tex, source lines 414--416. Cited by the retained statement or proof but not part of the collected spans; supplied only as a navigation pointer. The argument does not depend on its content.\par\medskip
\ifdefined\hypertarget\hypertarget{sourcepointer-3}{}\fi
\noindent\textit{Source reference.} Original lemma 0BW2: sources/stacks-discriminant/pinned-original-source.tex, source lines 1695--1699. Cited by the retained statement or proof but not part of the collected spans; supplied only as a navigation pointer. The argument does not depend on its content.\par\medskip
\ifdefined\hypertarget\hypertarget{sourcepointer-4}{}\fi
\noindent\textit{Source reference.} Original lemma 0BW1: sources/stacks-discriminant/pinned-original-source.tex, source lines 1657--1665. Cited by the retained statement or proof but not part of the collected spans; supplied only as a navigation pointer. The argument does not depend on its content.\par\medskip
\ifdefined\hypertarget\hypertarget{sourcepointer-5}{}\fi
\noindent\textit{Source reference.} Original lemma 0BVD: sources/stacks-discriminant/pinned-original-source.tex, source lines 509--513. Cited by the retained statement or proof but not part of the collected spans; supplied only as a navigation pointer. The argument does not depend on its content.\par\medskip
\ifdefined\hypertarget\hypertarget{sourcepointer-6}{}\fi
\noindent\textit{Source reference.} Original lemma 0BWL: sources/stacks-discriminant/pinned-original-source.tex, source lines 3188--3197. Cited by the retained statement or proof but not part of the collected spans; supplied only as a navigation pointer. The argument does not depend on its content.\par\medskip
\ifdefined\hypertarget\hypertarget{sourcepointer-7}{}\fi
\noindent\textit{Source reference.} Original lemma 0BWM: sources/stacks-discriminant/pinned-original-source.tex, source lines 3246--3255. Cited by the retained statement or proof but not part of the collected spans; supplied only as a navigation pointer. The argument does not depend on its content.\par\medskip
\begin{lemma}
\label{lemma-agree-dedekind}
Assume the Dedekind different is defined for $A \to B$.
Set $X = \Spec(A)$ and $Y = \Spec(B)$. The generalization of
Remark \ref{remark-different-generalization}
applies to the morphism $f : Y \to X$ if and only if
$1 \in \mathcal{L}_{B/A}$ (e.g., if $A$ is normal, see
Lemma \ref{lemma-dedekind-different-ideal}).
In this case $\mathfrak{D}_{B/A}$ is an ideal of $B$ and we have
$$
\mathfrak{D}_f = \widetilde{\mathfrak{D}_{B/A}}
$$
as coherent ideal sheaves on $Y$.
\end{lemma}

\begin{proof}
As the Dedekind different for $A \to B$ is defined we can apply
Lemma \ref{lemma-explain-condition} to see that
$Y \to X$ satisfies condition (1) of
Remark \ref{remark-different-generalization}.
Recall that there is a canonical isomorphism
$c : \mathcal{L}_{B/A} \to \omega_{B/A}$, see
Lemma \ref{lemma-dedekind-complementary-module}.
Let $K = Q(A)$ and $L = K \otimes_A B$ as above.
By construction the map $c$ fits into a commutative diagram
$$
\xymatrix{
\mathcal{L}_{B/A} \ar[r] \ar[d]_c & L \ar[d] \\
\omega_{B/A} \ar[r] & \Hom_K(L, K)
}
$$
where the right vertical arrow sends $x \in L$ to the map
$y \mapsto \text{Trace}_{L/K}(xy)$ and the lower horizontal
arrow is the base change map (\ref{equation-bc-dualizing}) for $\omega_{B/A}$.
We can factor the lower horizontal map as
$$
\omega_{B/A} = \Gamma(Y, \omega_{Y/X})
\to \Gamma(V, \omega_{V/X}) \to \Hom_K(L, K)
$$
Since all associated points of $\omega_{V/X}$
map to associated primes of $A$
(Lemma \ref{lemma-dualizing-associated-primes})
we see that the second map is injective.
The element $\tau_{V/X}$ maps to $\text{Trace}_{L/K}$ in
$\Hom_K(L, K)$ by the very definition of trace elements
(Definition \ref{definition-trace-element}).
Thus $\tau$ as in condition (2) of
Remark \ref{remark-different-generalization}
exists if and only if $1 \in \mathcal{L}_{B/A}$ and then
$\tau = c(1)$. In this case, by Lemma \ref{lemma-dedekind-different-ideal}
we see that $\mathfrak{D}_{B/A} \subset B$.
Finally, the agreement of $\mathfrak{D}_f$ with $\mathfrak{D}_{B/A}$
is immediate from the definitions and the fact $\tau = c(1)$ seen above.
\end{proof}

\end{document}
